%% Chapter 11 — Synthesis and Outlook
%%
%% This chapter is written fresh (no direct source paper).
%% It ties together all threads and makes explicit the postulate-reduction
%% argument across the whole monograph.

\chapter{Synthesis}
\label{ch:synthesis}

\section{Postulate Reduction: From Dozens to Two}
\label{sec:postulate-reduction}

% Conventional approaches (PK models, QSAR, pathway models, etc.) use
% many empirical postulates. This framework uses two axioms. Enumerate
% what is replaced.

% TODO

\section{The Computational--Biological Isomorphism}
\label{sec:isomorphism}

% Obs = Comp = Proc as the unifying equation of the monograph.
% Disease as computation failure.
% Therapy as computation restoration.

% TODO

\section{Connections to Established Theory}
\label{sec:connections}

% Poincare recurrence, Maslov, Kuramoto, Jaynes, Landauer, Shannon.
% What is new here vs. what is derived from these.

% TODO

\section{Open Problems}
\label{sec:open-problems}

% Genome-scale circuit topology requirement.
% Strong nonlinearity regime beyond first-order LP.
% Tissue-scale and whole-organ extension.
% Experimental apparatus for dual-pixel acquisition in clinical settings.

% TODO

\section{Falsifiability}
\label{sec:falsifiability}

% Single counterexample conditions that would refute the framework:
% - A drug whose effect cannot be expressed as edge conductance perturbation.
% - A disease with trivial loop holonomy.
% - A therapeutic outcome where variance does not drop.
% - Triple-observation divergence between optical, chromatographic, circuit
%   measurements on the same partition state.

% TODO
